\scanpage{134}
i.e. \(\unsure{1|\vert}\)
\[
 \begin{aligned} L_\omega&\simeq V\\ z&\longmapsto\Re z\end{aligned}
\]
est un isomorphisme d'espaces vectoriels complexes, quand on munit $V$ de la structure complexe $c_\omega$.\(\big]\)
% slanted
\margin{Mais il y a lieu d'associer à cette structure complexe sur $V$ \unsure{non pas|définie par|non par} $L_\omega$, mais $L_{\omega'}$, cf plus bas \ldots}

Dans le cas qui nous occupe, $\mathrm{rg}_{\mathbb R}V=2$ les \deleted{droites} sous-espaces vectoriels envisagés de $V_{\mathbb C}$ sont de \deleted{\unsure{rg}} rang $1$, \inserted{corang $1$,} et peuvent s'interpréter comme les points de
\[
 \Sigma=\Sigma_V=\mathbb P(\mathbb C),\quad\text{où}\quad P=\mathbb P(V) \tag*{(16)}
\]
est \deleted{l'}le \deleted{\unsure{esp}} schéma droite projective réelle associé à $V$. Pour que $\begin{aligned}L&\longrightarrow V\\ z&\longrightarrow\Re z\end{aligned}$ soit un isomorphisme, \deleted{i.e.} pour que $L$ corresponde à une structure complexe \deleted{de} sur $V$, il f. et s. que \deleted{\unsure{le} le point $z$ correspondant de $\Sigma$ soit} \inserted{$L\cap V=\{0\}$ i.e. $L$} n'est pas complexifié d'une \deleted{\unsure{ne provi|espèce|espace}} droite $L_0\subset V_{\deleted{\unsure{\mathbb R}}}$, i.e. que le point $z\in\Sigma$ qui correspond à $L$ \inserted{ne} soit pas dans \deleted{$\Sigma(\mathbb R)$}.
\[
 \Sigma_{\mathbb R}=P(\mathbb R)
\]
ou encore qu'on ait
\[
 z\in\Sigma\smallsetminus\Sigma_{\mathbb R}, \tag*{(17)}
\]

\deleted{\unsure{En résu|En résumé}} Ainsi, les structures complexes
% archive sheet number: 133
